\section{Autoencoders}
Autoencoders (AEs) are a class of neural networks trained in an unsupervised manner to learn efficient, compressed representations of data. 

\paragraph{The Core Idea}
Conceptually, autoencoders are very similar to Principal Component Analysis (PCA) in that they aim to learn a lower-dimensional representation of the input data, known as the \textbf{latent space}. 

While PCA is restricted to discovering strictly linear transformations, autoencoders leverage non-linear activation functions, allowing them to model much more complex, non-linear data manifolds. In fact, if an autoencoder uses only linear activations and is trained with a mean squared error loss, it learns to span the exact same principal subspace as PCA.

This lower-dimensional representation is learned by training the network to map the input data $x$ to a compressed space $z$ via an \textit{encoder}, and then reconstruct the original data as accurately as possible from that representation via a dual network called the \textit{decoder}.

\begin{comment}
    #TODO: add figure of autoencoder architecture
\end{comment}

To train an autoencoder, we minimize the reconstruction loss between the input and the output of the network. 
Formally, given an input $x$, the encoder $f_\theta$ maps it to a latent representation $z = f_\theta(x)$, and the decoder $g_\phi$ reconstructs the input as $\hat{x} = g_\phi(z)$. 

The training objective is to minimize the reconstruction loss, which is typically defined using the Mean Squared Error (MSE):
\begin{equation}
    \mathcal{L}(x, \hat{x}) = ||x - \hat{x}||^2
\end{equation}

\paragraph{Implementation}
Autoencoders give us a lot of flexibility in terms of architecture and design choices. In fact, while the idea is fixed, 
the implementation can vary significantly depending on the specific use case and data type.


\paragraph{Why AEs?}
Autoencoders have been shown to be very effective feature extractors, and they can be used for a variety of tasks, including dimensionality reduction, anomaly detection, and generative modeling.
To do so, the decoder is usually discarded after training. The encoder can then be used to map new data points to the learned latent space.


\subsection{Variational Autoencoders}
From the idea of autoencoders, we can derive a more powerful generative model called the \textbf{Variational Autoencoder} (VAE).
This method belongs to the class of \textbf{explicit density estimation} generative models, which directly learn the probability distribution of the data.


\paragraph{Idea}
VAEs define a density function $p_\theta(x)$ over the data space, parameterized by a neural network with parameters $\theta$. 

\begin{equation}
    \label{eq:VAEs}
    p_\theta(x) = \int p(z)\, p_\theta(x \mid z)\, dz,
\end{equation}

where $p(z)$ is a prior distribution over the latent variable $z$ (typically a standard Gaussian). 

We would like to optimize the model with respect to $\theta$, but this is intractable. 
The problem is that the integral in Equation \ref{eq:VAEs} cannot be computed in closed form when $p_\theta(x \mid z)$ is parameterized by a neural network. 
You cannot integrate over all possible latent values $z$.

A natural next thought is to use the \textit{posterior} distribution $p_\theta(z \mid x)$. Applying Bayes' rule, we can write the posterior as:
\begin{equation}
    p_\theta(z \mid x) = \frac{p_\theta(x \mid z)\, p(z)}{p_\theta(x)}.
\end{equation}

But again, the denominator $p_\theta(x)$ is intractable, so we cannot compute the posterior either.

Since both the marginal likelihood and the posterior are intractable, we cannot directly optimize the model with respect to $\theta$.
We can, however, introduce a new distribution $q_\phi(z \mid x)$, parameterized by a second neural network with parameters $\phi$, which will serve as an approximation to the true posterior $p_\theta(z \mid x)$.

\paragraph{A tractable solution}
Given this approximation, we can rewrite the marginal likelihood as:
\begin{equation}
    \log p_\theta(x) = E_z \left[ \log p_\theta(x \mid z) \right] - D_{KL}\left( q_\phi(z \mid x) || p(z) \right) + D_{KL}\left( q_\phi(z \mid x) || p_\theta(z \mid x) \right),
\end{equation}

where $D_{KL}$ is the Kullback-Leibler divergence, which measures the difference between two probability distributions.

\begin{itemize}
    \item The first term is the expected log-likelihood of the data given the latent variable $z$, which encourages the model to reconstruct the input data accurately.
    It can be computed by sampling $z$ from the approximate posterior $q_\phi(z \mid x)$.
    \item The second term is the KL divergence between the approximate posterior and the prior, which encourages the approximate posterior to be close to the prior distribution.
    \item The third term cannot be computed directly, but it is always non-negative.
\end{itemize}

\paragraph{Evidence Lower Bound (ELBO)}

To overcome the intractability of the marginal likelihood, VAEs maximize a tractable lower bound on the log-likelihood, known as the \textbf{Evidence Lower Bound} (ELBO):

\begin{equation}
    \mathbb{E}_{z} \left[\log p_\theta(x|z)\right] - D_{KL} \left(q_\phi(z|x)\,\|\,p(z)\right).
\end{equation}

The ELBO consists of two terms:

\begin{itemize}
    \item The first term is the \textit{reconstruction term}. It measures how well the decoder $p_\theta(x|z)$ can reconstruct the observed data $x$ from a latent representation sampled from the approximate posterior $q_\phi(z|x)$. Maximizing this term encourages accurate reconstructions.

    \item The second term is the \textit{regularization term}. It is the Kullback--Leibler (KL) divergence between the approximate posterior $q_\phi(z|x)$ and the prior distribution $p(z)$. Minimizing this term encourages the latent representations to remain close to the prior, resulting in a smooth and well-structured latent space.
\end{itemize}

This loss function allows VAEs to learn a meaningful latent space, where similar inputs are mapped to nearby points in the latent space.

\paragraph{Inference Time}